One of the desired outcomes for this work is formulation of \gls{DNP} group parameters for \glspl{MSR} which can be applied generally.
If the \gls{DNP} group parameters can be applied generally to different systems with similar residence times, then the amount of computational work required is drastically reduced.
However, it is necessary that sufficient evidence is provided demonstrating that the generated parameters are identical for different \glspl{MSR}.

Specifically, the \gls{MSRE} and the \gls{MSFR} will be simulated using the spatially resolved model \cite{robertson1965msre, brovchenko2019neutronic}.
The models will be adjusted such that the in-core and ex-core residence times are identical for each.
The result will be a set of group parameters for two distinct reactors but with the same residence time tabulations.
Instead of calculating the neutron flux, both reactors will be given an incident neutron energy spectra for irradiating the sample.
The impact of the incident neutron energy spectra is discussed in Section \ref{sec:indidentspectra}, so differences caused by the spectra will not be studied in this section.
Additionally, the fuel salt and fuel temperature in the \gls{MSFR} will be replaced with those from the \gls{MSRE}.

The two different reactor simulations will have all of the same input parameters aside from the mesh and thermal-hydraulic parameters.
If the \gls{DNP} group parameters have negligible differences, then the geometry and thermal-hydraulic parameters are not necessary when generating the group parameters.
Instead, all that is needed are the other input parameters, meaining the work does not need to be replicated for every different reactor geometry and thermal-hydraulic configuration.
If the \gls{DNP} group parameters have non-negligible differences, then each \gls{MSR} needs to have \gls{DNP} group parameters generated for its own specific geometry and flow conditions.
Further investigation may also be warranted to determine if there are one or two parameters which could be used to capture these differences, thus reducing the total amount of work necessary to model all \gls{MSR} designs.
